\chapter{Einleitung}
\sloppy
Die Fähigkeit von IT-Infrastrukturen, sich dynamisch an variierende Anforderungen anzupassen, ist eine zentrale Anforderung zeitgemäßer Softwarearchitektur (vgl. Bharadwaj et al. 2013, S. 471 f.). Unternehmerische Agilität, verstanden als die Kapazität einer Organisation, schnell und effektiv auf technische wie fachliche Veränderungen zu reagieren, setzt voraus, dass Systeme bedarfsgerecht skaliert werden können, ohne dauerhaft physische Ressourcen binden zu müssen (vgl. Tallon/Pinsonneault 2011, S. 463 f.). Gleichzeitig gewinnt die digitale Souveränität an Bedeutung, die den selbstbestimmten Umgang einer Organisation mit ihren Infrastrukturentscheidungen, Daten und Technologieabhängigkeiten bezeichnet (vgl. Couture/Toupin 2019, S. 2 ff.). Während klassische Systeme oft starre Strukturen aufweisen, erfordern Big-Data-Szenarien heuzutage Architekturen, die Souverinität und Agilität inhärent abbilden (vgl. Marz/Warren 2015, S. 4 ff.).

Als empirische Grundlage dient der NYC Taxi Trip Data-Datensatz, der umfangreiche Informationen zu Taxifahrten in New York City enthält. Anhand dieses Datensatzes werden Datenoperationen von der Berechnung deskriptiver Kennzahlen über explorative Visualisierungen bis hin zum Training von Vorhersagemodellen durchgeführt, um die Leistungsfähigkeit der Python-basierten, containerisierten Systemarchitektur unter realitätsnahen Lastbedingungen zu demonstrieren. Im Vordergrund stehen dabei nicht die domänenspezifischen Erkenntnisse aus den Verkehrsdaten selbst, sondern die systemarchitektonischen Herausforderungen, eine robuste und zugleich flexible Umgebung für Softwareentwicklung und Betrieb zu schaffen. Der Einsatz von Docker-Containern gewährleistet dabei Portabilität, elastische Skalierbarkeit und nahtlose Plattformunabhängigkeit (vgl. Merkel 2014, S. 2 ff.).

Ziel dieser Arbeit ist die Konzeption und prototypische Implementierung einer modularisierten Infrastruktur, die durch konsequente Diensttrennung und Orchestrierung eine skalierbare Datenverarbeitungspipeline realisiert, vom Datenimport über die Aufbereitung bis zur Bereitstellung via API sowie der finalen Visualisierung und Vorhersage (vgl. Newman 2015, S. 1 ff.).